% Propietario conservado: /Users/ruben/Documents/New project/output/EXTREMA_MEDIA_RAZON_RECIPROCIDAD_ES_EN_20260918/ARTICULO/ES/source/libro/reciprocidad_apertura.tex
\HMTDivision{\(K\;\longleftrightarrow\;\Lambda_t\;\longleftrightarrow\;\Lambda_t^*\)}{Geometría recuperable y reciprocidad del acoplamiento}{Deformación positiva, dualidad reticular, estructura polar y recuperación orientada de la acción.}
\label{x_reciprocidad:rec:apertura}

La función de la constante holográfica de acoplamiento aritmético--geométrico se precisa al estudiar las transformaciones que su registro determina. Su lectura posicional conserva doce componentes ordenadas; su realización incidencial distingue una dirección; el archivo bilateral permite reconstruirla. La composición de estas operaciones conduce a una familia de geometrías cuyo volumen fundamental permanece constante y cuya inversión produce la red dual. El registro interviene, por tanto, tanto en la constitución de la geometría como en su recuperación.

Este desarrollo mantiene el orden generativo establecido en el artículo. Las evaluaciones aditiva y multiplicativa de APP, la orientación del TRIT y las operaciones de selección, transporte y memoria del TPK producen el estado enriquecido y su estructura discreta conjunta del continuo. Las coordenadas regionales y la incidencia excepcional son lecturas de ese estado. Las realizaciones vectoriales, reticulares e hiperbólicas que siguen reciben esas salidas con sus cartas y sus orientaciones. La genealogía se conserva al cambiar de representación.

La distinción entre escala y forma adquiere aquí una expresión exacta. Una deformación puede dilatar determinadas direcciones y contraer otras mientras conserva el producto total de sus factores. En la elipse de acción, el producto de los semiejes conserva la sección de acción. En los doce radios determinados por \(K\), la compensación conserva el covolumen reticular. En el operador de masa, una descomposición separa el factor uniforme de la transformación relativa de determinante uno. Cada conservación tiene un objeto y una demostración propios; los lectores permiten relacionarlas sin confundir sus dominios.

La reciprocidad radial articula dos construcciones complementarias. En la carta del doble círculo espectral, invertir el radio equivale exactamente a reflejar la coordenada compleja respecto de la recta de parte real \(1/2\). En la carta modular de la elipse, intercambiar semiejes invierte el radio y cambia el signo de la coordenada de Klein. Ambas correspondencias admiten fórmulas inversas. El registro de orientación distingue la lectura conjugada de una modificación de la sección física de acción.

Al sustituir los lectores angulares por sus expresiones completas aparece una cancelación aritmética específica: el término lineal de la sección previa de acción se compensa con el término correspondiente del contraángulo. La continuación regional permanece íntegra. El contraste angular, el radio y la coordenada hiperbólica quedan entonces expresados mediante las mismas publicaciones que intervienen en \(\alpha\), con la constante de acoplamiento dentro de su clausura dodecafásica. Recíprocamente, la coordenada radial orientada recupera la sección de retorno y permite reescribir la expresión electrónica completa.

La realización espectral distingue a su vez dos clases de información. La dualidad de redes impone una ecuación funcional; el covolumen fija los polos y sus residuos; la deformación modifica la parte entera de la función completada. El archivo recupera la dirección que determina esa variación y proporciona cotas explícitas para el error de reconstrucción. La conservación y la reversibilidad se expresan así mediante operaciones que abarcan el registro, la geometría y sus lecturas analíticas.

\HMTMathematicalPaper
